\documentclass[11pt,a4paper]{extarticle}

\usepackage[T2A]{fontenc}
\usepackage[utf8]{inputenc}
\usepackage[russian]{babel}
\usepackage{geometry}
\usepackage{graphicx}
\usepackage{booktabs}
\usepackage{amsmath,amssymb}
\usepackage{hyperref}
\usepackage{float}
\usepackage{enumitem}
\usepackage{titlesec}
\usepackage{microtype}
\usepackage{xcolor}

\geometry{left=19mm,right=19mm,top=17mm,bottom=18mm}
\setlength{\parindent}{1cm}
\setlength{\parskip}{0.25em}
\definecolor{linkblue}{RGB}{35,88,150}
\hypersetup{colorlinks=true, linkcolor=linkblue, urlcolor=linkblue}
\titleformat{\section}{\Large\bfseries}{\thesection}{0.6em}{}
\titleformat{\subsection}{\large\bfseries}{\thesubsection}{0.5em}{}
\setlist[itemize]{topsep=2pt,itemsep=1pt,leftmargin=18pt}

\newcommand{\maybefig}[3]{%
  \IfFileExists{../output/#1}{%
    \begin{figure}[H]\centering
    \includegraphics[width=#2\textwidth]{../output/#1}
    \caption{#3}
    \end{figure}
  }{}
}

\newcommand{\maybepair}[4]{%
  \IfFileExists{../output/#1}{%
    \IfFileExists{../output/#2}{%
      \begin{figure}[H]\centering
      \begin{minipage}{0.49\textwidth}\centering
      \includegraphics[width=\linewidth]{../output/#1}
      \end{minipage}\hfill
      \begin{minipage}{0.49\textwidth}\centering
      \includegraphics[width=\linewidth]{../output/#2}
      \end{minipage}
      \caption{#3 (слева) и #4 (справа).}
      \end{figure}
    }{}
  }{}
}

\begin{document}

\begin{center}
{\LARGE\bfseries ML Start 2026: домашнее задание \#3}\\[2mm]
{\Large Тема 24: глубина vs ширина MLP}\\[1mm]
{\normalsize Сорочан Дмитрий}
\end{center}

\section{Постановка задачи}

\hspace*{11mm}Цель работы --- исследовать, как глубина и ширина многослойного перцептрона влияют на качество. Используются оба корпуса из темы 24 --- MNIST и Fashion-MNIST --- и отдельный блок с синтетическими функциями. Для классификации вход имеет 784 пикселя, выход --- 10 классов.

Базовая архитектура намеренно простая:
$$
  x 
  \to \mathrm{Flatten} 
  \to (\mathrm{Linear}
  \to\mathrm{ReLU})^d 
  \to \mathrm{Linear}(10).
$$
Фиксированы Adam, learning rate и weight decay. Это позволяет менять именно $d$ и $w$, а не одновременно архитектуру, scheduler, BatchNorm и residual connections.

Для одинаковой ширины число параметров зависит от глубины:
$$
    P(d, w)
    = (784 + 1) \cdot w + (d - 1) \cdot (w + 1) \cdot w + (w + 1) \cdot 10.
$$
Поэтому кроме обычного sweep проведен parameter-matched эксперимент, где shallow и deep сети сравниваются при близком числе параметров.

\subsection{Теоретическая мотивация}

\hspace*{11mm}Теорема об универсальной аппроксимации говорит, что достаточно широкая сеть с одним скрытым слоем и подходящей нелинейностью может приближать непрерывные функции на компакте. Но теорема не говорит, какая ширина нужна, как найти веса и насколько хорошо модель будет обобщать. Глубина может быть параметрически выгоднее, но очень глубокую plain MLP тяжелее оптимизировать. Работа проверяет именно этот компромисс.

\section{Данные, EDA и признаки}

\hspace*{11mm}MNIST и Fashion-MNIST имеют одинаковый формат: $60\,000$ train, $10\,000$ official test, изображения $28\times28$, 10 классов. Пропусков и категориальных признаков нет. Fashion-MNIST визуально сложнее из-за похожих классов верхней одежды.

\maybepair{eda_class_balance.png}{eda_engineered_features_correlations.png}{баланс классов}{корреляции engineered features}
\maybepair{eda_mnist_engineered_features_by_class.png}{eda_fashion_engineered_features_by_class.png}{MNIST: engineered features по классам}{Fashion-MNIST: engineered features по классам}

Для классических моделей к 784 пикселям добавлены пять image-level признаков: средняя яркость, стандартное отклонение пикселей, доля активных пикселей и координаты центра массы. Коррелированные признаки не удаляются автоматически: для деревьев корреляция не нарушает корректность обучения, но importance может распределяться между похожими признаками.

В EDA дополнительно строятся распределения каждого engineered feature по 10 классам отдельно для MNIST и Fashion-MNIST. Это позволяет проверить, несут ли глобальная заполненность, контраст и центр массы class-conditional сигнал до обучения классических моделей.

Основная метрика --- \textbf{balanced accuracy}; дополнительно считаются macro-F1, macro-precision и macro-recall. Хотя оба корпуса почти сбалансированы, единая метрика делает все сравнения согласованными.

\section{Глубина и ширина MLP}

\hspace*{11mm}Основная прямоугольная сетка:
$$
  d
  \in \{1, 2, 3, 5, 10, 20, 50, 100\},\qquad
  w
  \in \{8, 16, 32, 64, 128, 256, 512, 1024, 2048\}.
$$
На каждой точке фиксированы train subset, validation, число эпох, optimizer и regularization. Early stopping в sweep не используется, чтобы разные архитектуры получили одинаковый training budget.

\maybepair{mnist_val_balanced_accuracy_heatmap.png}{fashion_val_balanced_accuracy_heatmap.png}{MNIST: width $\times$ depth}{Fashion-MNIST: width $\times$ depth}

Широкие сети входят в тот же основной sweep, поэтому отдельного второго прохода нет. Для сочетаний свыше 20 млн параметров обучение пропускается до создания модели: например, $d=100,w=2048$ потребовал бы сотни миллионов параметров и многогигабайтные состояния Adam. Такие точки проверяли бы прежде всего доступную память, а не архитектурный эффект.

\subsection{Learning curves, регуляризация и CV}

\hspace*{11mm}Weight decay $10^{-4}$ используется одинаково для всех MLP. На learning curves train и validation loss строятся на одном графике; отдельно сравнивается balanced accuracy. Это позволяет отличить переобучение от ситуации, когда deep-сеть просто плохо оптимизируется.

\maybepair{mnist_learning_curve_loss.png}{fashion_learning_curve_loss.png}{MNIST: train/validation loss}{Fashion-MNIST: train/validation loss}

Большой sweep выполняется на одном фиксированном stratified hold-out split. После выбора архитектура дополнительно проверяется 5-fold StratifiedKFold на train subset. Это не используется как новый поиск hyperparameters, а служит проверкой устойчивости выбранной конфигурации.

Official test не участвует ни в sweep, ни в CV. После выбора сеть переобучается на всей official train-части и test оценивается один раз.

\section{Классические модели}

\hspace*{11mm}Проверяются GaussianNB, Logistic Regression и Random Forest. Внешний validation отделяется заранее; GridSearchCV работает только на train через 5-fold StratifiedKFold. StandardScaler для Logistic Regression находится внутри Pipeline.

Для Random Forest не задается искусственный лимит глубины. GridSearchCV выполняется последовательно (\texttt{SEARCH\_N\_JOBS=1}), а параллелизм остается только внутри самого леса. Это избегает nested joblib-parallelism и лишних пиков памяти.

\maybefig{mnist_random_forest_feature_importance.png}{0.72}{MNIST: Random Forest feature importance.}

Feature importance имеет реальные имена: \texttt{pixel\_row\_col} и пять engineered features.

\section{Confusion matrix и анализ ошибок}

\hspace*{11mm}Для лучших моделей строятся normalized confusion matrix и таблица наиболее частых направленных ошибок \texttt{true -> predicted}. На Fashion-MNIST особенно важны пары Shirt/T-shirt/top/Pullover/Coat; это согласуется с EDA. На MNIST ошибки менее выражены и распределены между визуально похожими написаниями цифр.

\maybepair{mnist_best_classical_confusion_matrix.png}{fashion_mlp_confusion_matrix.png}{MNIST: лучшая классика}{Fashion-MNIST: MLP}

\section{Синтетические функции}

\hspace*{11mm}Выбраны четыре структуры: почти линейная функция, полином, периодическая функция и Heaviside. Теперь одна и та же сетка $d\times w$ прогоняется \textbf{по всем четырем функциям}, а не только по periodic. Это важно: оптимальная архитектура может зависеть от гладкости и локальной структуры целевой зависимости.

\maybefig{synthetic_functions.png}{0.70}{Выбранные синтетические функции.}
\maybefig{synthetic_mse_heatmaps_all.png}{0.88}{MSE heatmap width $\times$ depth для всех четырех функций.}

В старой версии конфигурация $d=20,w=32$ иногда превращалась почти в горизонтальную прямую. Проверка показала, что это был реальный optimization collapse, а не баг визуализации: стандартная инициализация вместе с 20 слоями Tanh почти уничтожала вариативность сигнала. Для всех конфигураций одинаково введена Xavier initialization с gain для Tanh и $LR=3\cdot10^{-3}$. В контрольном запуске periodic для $d=20,w=32$ получила MSE около $1.6\cdot10^{-3}$ и перестала быть константой. Большой synthetic sweep запускается на CPU: для 1D-данных GPU не нужен, а это исключает накопление MPS allocator cache при последовательном создании десятков моделей.

Heaviside остается принципиально особым примером: конечная MLP с непрерывными активациями непрерывна, поэтому точный скачок в sup-норме она не воспроизводит, хотя MSE может быть малым за счет узкой переходной области.

\section{Blending и корреляции моделей}

\hspace*{11mm}Проверяются equal-weight и metric-weighted soft voting. Metric-based численные веса пропорциональны CV balanced accuracy на train. Однако MLP-архитектура до этого выбиралась по hold-out validation, поэтому весь blend-score трактуется как exploratory, а не как независимая test-оценка ансамбля. Для строгой финальной оценки потребовался бы отдельный untouched split для blending.

Перед blending строятся две таблицы: корреляция развернутых probability predictions и корреляция бинарных индикаторов ошибок. Коррелировать номера классов нельзя --- это категориальные метки, а не числовая шкала.

\maybepair{mnist_prediction_correlations.png}{fashion_prediction_correlations.png}{MNIST: корреляции predictions/errors}{Fashion-MNIST: корреляции predictions/errors}
\maybefig{fashion_mnist_blending_metrics.png}{0.80}{Fashion-MNIST: одиночные модели и варианты blending.}

В сохраненном прогоне probabilities LR/RF/MLP высоки, примерно $0.94$--$0.95$, а корреляции индикаторов ошибок умеренные. Поэтому большой прирост от простого усреднения не ожидается. На MNIST лучшая одиночная MLP осталась сильнее blend; на Fashion-MNIST equal blend улучшил balanced accuracy примерно с $0.8782$ до $0.8804$, metric-weighted --- до $0.8805$. Отсутствие большого прироста здесь является содержательным отрицательным результатом, а не ошибкой кода.

\section{Корректность, leakage и воспроизводимость}

\begin{itemize}
\item единый \texttt{SEED=143} для split, CV, NumPy, Python и PyTorch;
\item все сравниваемые корпуса используют одинаковый hold-out validation protocol;
\item official test не участвует в выборе архитектуры, hyperparameters или blend weights;
\item StandardScaler обучается только внутри CV Pipeline;
\item GridSearchCV видит только train, внешний validation остается контролем;
\item metric-based blend вычисляет численные веса из train-CV scores; поскольку MLP-кандидат выбран по validation, сами blend-метрики явно помечены как exploratory;
\item перед blending проверяется идентичность порядка \texttt{y\_true} во всех NPZ;
\item synthetic initialization и learning rate одинаковы для всех depth/width;
\item результаты сохраняются в \texttt{/mlStart/hws/output} независимо от cwd kernel.
\end{itemize}

\section{Чек-лист задания и итог}

\hspace*{11mm}Выполнены EDA, feature engineering, три классических алгоритма, GridSearchCV и 5-fold CV, MLP с learning curves и регуляризацией, сравнение качества/времени, confusion matrix и error analysis. Для темы 24 есть MNIST, Fashion-MNIST, synthetic functions, единый sweep с ширинами 8--2048, heatmap, parameter-matched shallow/deep и MSE heatmap по всем синтетическим функциям. Бонусы закрыты feature importance, blending, correlation analysis и GitHub-ready оформлением.

\textbf{Главный вывод:} глубину и ширину нельзя оценивать отдельно от training budget и числа параметров. Ширина обычно дает насыщение, тогда как чрезмерная глубина plain MLP может ухудшать оптимизацию. Но и этот эффект надо отделять от технических артефактов вроде неудачной инициализации --- именно поэтому synthetic block использует одинаковую Xavier initialization и отдельные learning curves.

\IfFileExists{../output/generated_results.tex}{\input{../output/generated_results.tex}}{}

\end{document}
